\scanpage{204}
\mspage{78}
d'autre part tautologiquement
\diagram{\begin{tikzcd}[ampersand replacement=\&]
 a\arrow[r,Rightarrow]\&b\arrow[r,Rightarrow]\arrow[d,Rightarrow]\&b'\arrow[d,Rightarrow]\\
 \&c\arrow[r,Rightarrow]\&c'
\end{tikzcd}}
et il reste à prouver $c'\Longrightarrow a$, ce qu'on fera plus bas.

\unsure{Prouvons|Revenons} donc 6.1. Dans (6.1.1) l'égalité $\widetilde C_0=\widetilde{C^{\unsure{\pi|R}}}$ est tautologique. Donc tout revient à prouver
\[
 C\subset\widetilde{C^\pi}.
\]
une difficulté ici, c'est qu'on ne sait \uline{pas} que $\widetilde{C^\pi}$ soit stable par $\underrightarrow{\lim}$ filtrantes --- on le sait seulement pour $C$.

Soit donc
\[
 A\lhook\joinrel\longrightarrow B
\]
une flèche dans $C$, prouvons qu'elle est dans $\widetilde{C^\pi}$. Pour ceci, on \deleted{\unsure{pratique}} va construire, par induction transfinie, une suite transfinie \inserted{\uuline{\unsure{strictmt|strictemt}} croissante} de sous-objets $A_\alpha$ de $B$ contenant $A=A_0$, \deleted{et telle qu'on ait}
\diagram{\begin{tikzcd}[ampersand replacement=\&]
 A\arrow[r,hook,"{i_\alpha}"]\arrow[d,equal]\&A_\alpha\arrow[r,hook,"{\unsure{i_{\beta\alpha}|j_{\beta\alpha}}}"]\&\deleted{\unsure{B|A}}A_\beta\arrow[r,"{j_\beta}"]\&B\\
 A_0\&\&\&
\end{tikzcd}}
telle que l'on ait
\[
 i_\alpha\in\widetilde{C^{\pi}_{\unsure{4|\ast|\omega}}},\quad
 \inserted{\text{plus gén. les }i_{\beta,\alpha}\in\widetilde{C^\pi},}\quad j_\alpha\in C
\]
et que l'on ait
\[
 A_\alpha=\underset{\beta<\alpha}{\underrightarrow{\lim}}A_\beta
 \qquad\text{si }\deleted{\unsure{\beta}}\alpha\text{ ordinal limite.}
 \tag*{\unsure{(6.2.)|(6.2)}}
\]
% archive sheet number: 203
